% Vietnamese translation for OI Public Library
% Source-File: IOI中国国家候选队论文/2006/唐文斌/唐文斌.ppt
\documentclass[11pt]{article}
\usepackage{oipl-vn}

\title{Tư tưởng điều chỉnh trong thi tin học\\{\large Ghi chú theo slide}}
\author{Tang Wenbin}
\date{2006}

\begin{document}
\maketitle

\origtitle{Adjustment thinking in informatics contests}
\sourcepath{IOI China candidate team papers / 2006 / Tang Wenbin / Tang Wenbin.ppt}

\section{Phạm vi}

Bộ slide đi kèm bài viết về cách bắt đầu từ một phương án thô rồi điều chỉnh
dần tới phương án hợp lệ hoặc tốt hơn. Ghi chú này dựa trên bản bài viết đã
dịch.

\section{Ý tưởng chung}

Trong một số bài, xây trực tiếp nghiệm tối ưu quá khó. Ta có thể tạo một cấu
hình ban đầu dễ có, sau đó áp dụng các phép sửa cục bộ. Mỗi phép sửa phải có
tác dụng rõ ràng: giảm số lỗi, cải thiện hàm mục tiêu, hoặc đưa cấu hình tới
một miền dễ phân tích hơn.

\section{Dựng nghiệm}

Ở các bài dựng nghiệm, điều chỉnh thường đi kèm một đại lượng đơn điệu. Ví dụ
trong bài viễn thông, bắt đầu bằng cách đặt toàn bộ một phía gửi và phía kia
nhận. Mỗi lần gặp một cổng nhận chưa hợp lệ, đổi trạng thái của nó và cập
nhật cổng liên quan. Số lỗi giảm theo một hướng nên thuật toán dừng và nghiệm
cuối cùng hợp lệ.

\section{Bài mở và mô phỏng}

Với bài không yêu cầu tối ưu tuyệt đối, điều chỉnh có thể được dùng như một
chiến lược cải tiến nghiệm: thử thay đổi nhỏ, đánh giá điểm số, rồi giữ lại
thay đổi tốt. Mô phỏng tôi luyện mở rộng ý tưởng này bằng cách đôi khi chấp
nhận bước kém hơn để thoát khỏi cực trị cục bộ.

\section{Điều cần kiểm soát}

Một thuật toán điều chỉnh cần ba phần: cách sinh trạng thái ban đầu, tập phép
sửa, và tiêu chí bảo đảm tiến triển. Nếu thiếu tiêu chí tiến triển, chương
trình có thể dao động; nếu phép sửa quá nghèo, nó có thể dừng ở cấu hình chưa
đạt yêu cầu.

\end{document}
